\begin{scanpage}{0074}{57}
Assuming $\phi=-\frac\pi2$ and $\theta=\pi-\gamma$ $\left(\frac\pi2<\theta<\pi\right)$ in the resulting expressions, let us find the field scattered by the disk in the direction toward its source. With E-polarization of the incident wave, it equals
\[\left.\begin{aligned}E_\phi=-H_\theta&=-\frac{iaE_{0x}}2\frac{\cos\theta}{\sin\theta}J_1(2ka\sin\theta)\frac{e^{ikR}}R,\\E_\theta&=H_\phi=0,\end{aligned}\right\}\tag{10.06}\]
and with H-polarization
\[\left.\begin{aligned}E_\theta=H_\phi&=\frac{iaH_{0x}}2\frac{\cos\theta}{\sin\theta}J_1(2ka\sin\theta)\frac{e^{ikR}}R,\\E_\phi&=H_\theta=0.\end{aligned}\right\}\tag{10.07}\]
Using the asymptotic expressions for the Bessel functions, one is able to show that when $R\gg ka^2$ and $ka\sqrt{\lambda^2+\mu^2}\gg1$ the fringing field is radiated from a luminous region on the disk. In the case when $\sqrt{\lambda^2+\mu^2}\to0$, the luminous region is increased and in the limit (when $\lambda=\mu=0$) the entire surface of the disk starts to ``shine''.
\sectiontitle{11}{The Field Radiated by the Nonuniform Part of the Current}
Let us calculate the field in the nonuniform part of the current
\[\mathbf j^1(\rho,\psi)=\mathbf J(\rho,\psi)e^{ik\rho\sin\gamma\sin\psi}.\tag{11.01}\]
Its corresponding vector potential
\[\mathbf A=\frac ac\frac{e^{ikR}}R\int_0^a d\rho\int_0^{2\pi}\mathbf J(\rho,\psi)e^{-ik\rho\sqrt{\lambda^2+\mu^2}\cos(\psi-\delta)}\,d\psi\tag{11.02}\]
by means of the stationary phase method is transformed with $ka\sqrt{\lambda^2+\mu^2}\gg1$ to the form
\[\begin{aligned}\mathbf A&=\frac ac\sqrt{\frac{2\pi}{ka\sqrt{\lambda^2+\mu^2}}}e^{i\frac\pi4}\left[\int_0^a\mathbf J(\rho,\psi_1)e^{-ik\rho\sqrt{\lambda^2+\mu^2}}\,d\rho-\right.\\&\qquad\left.-i\int_0^a\mathbf J(\rho,\psi_2)e^{ik\rho\sqrt{\lambda^2+\mu^2}}\,d\rho\right]\frac{e^{ikR}}R.\end{aligned}\tag{11.03}\]
Here
\end{scanpage}
